\section{Background}

\subsection{Offline vs Online Calibration}

A scorer is \textit{offline-calibrated} when its weights are fit on data
that includes or post-dates the period it is applied to; it is
\textit{online-calibrated} when its weights are fit using only information
available before the prediction time. The distinction matters because
deployable procedures cannot peek into the future. Calibration of
scoring models is itself a well-studied problem: modern classifiers
are frequently miscalibrated~\cite{guo2017calibration}, and a range of
post-hoc remedies, including Platt scaling~\cite{platt1999probabilistic},
isotonic-regression score transformation~\cite{zadrozny2002transforming},
and the broader empirical comparison of probability-producing
learners~\cite{niculescu2005predicting}, adjust raw scores toward
well-calibrated probabilities. Those methods fit a static correction on
a held-out set; the present work instead re-fits the scorer's
combination weights as the operational state evolves. The
temporal-stability study's H3
ablation computed per-window weights by grid search
on the calibration seeds' \textit{same-window} state, an offline-peek
procedure that establishes an upper bound on what real-time recalibration
could achieve, but is not itself deployable.

\subsection{Rolling-History Cross-Validation}

In rolling-history calibration with a one-window lag, at each scoring
window $w$ the weights used are fit on window $w-1$ only of a held-out
calibration set; window $w$ data is never inspected at training time.
This pattern is standard in time-series cross-validation and is the
weakest form of \textit{strictly online} learning: it uses only past
observations and is therefore a valid stand-in for a real-time
operational deployment. Re-fitting weights as the underlying
distribution shifts is an instance of the concept-drift adaptation
problem~\cite{gama2014survey}; the lag-1 grid search studied here is a
deliberately minimal, auditable point in that design space.

\subsection{Why High-Capacity Cells Matter Most}

The capacity-decay study~\cite{paper6capacity} documented that HygienePrio-full's
absolute W6 P@50 collapses at $K \in \{100, 200\}$ even though its
per-pair relative advantage over EPSS-only is preserved. The natural
follow-up question, can online recalibration rescue HygienePrio at
high capacity?, is operationally important because organisations
with large patching teams sit precisely in the cells where the
fixed-weight procedure collapses.

\subsection{Policy Mandates}

CISA BOD~22-01~\cite{cisa2021bod2201} and 23-01~\cite{cisa2022bod2301}
together with NIST SP~800-40 Rev.~4~\cite{nist2022sp80040r4} define the
operational scheduling regime under which any recalibration procedure
must work. Stakeholder-specific prioritization schemes such as
SSVC~\cite{spring2021ssvc} likewise assume a deployable, auditable
decision procedure rather than a future-peeking one. The procedure must be deployable, must not depend on
future-window data, and must be auditable, conditions that the
rolling-history grid search satisfies and offline-peek does not.
